% 05_defi_experiment.tex --- DeFi Experiment (long-form analogue of
% ICICPE paper Section V). Narrates forecaster training, headline H1
% test, and the hysteresis-sweep ablation in roughly 1200 words. All
% numeric entries are pulled either from the macros defined in
% sections/09_results.tex (auto-filled from results/tables/*.csv) or
% from forecaster/trained_models/metrics.json directly.

\section{DeFi Experiment}
\label{sec:defi-experiment}

This section narrates the empirical evaluation of the architecture,
forecaster, and allocator developed in
\cref{sec:methodology,sec:dual-branch,sec:mcdm,sec:decision} on the
Aave V3 / Compound V3 USDC supply panel of \cref{sec:data}. It is the
long-form companion to the headline summary tables of
\cref{sec:results}: where \cref{sec:results} reports the binding
numbers in tabular form, this section unpacks the training and
validation diagnostics, the bootstrap, the per-quarter regime
breakdown, and the hysteresis-sweep ablation as a connected narrative.
The headline pre-registered test
$H_1{:}\,\Delta\textsc{Sharpe}\ge 0.2$ from \cref{sec:introduction}
is evaluated on a binding four-month held-out fold, and the
publication protocol of \cref{sec:contributions} commits to reporting
the result whether or not $H_1$ is rejected. As we shall see, the
test window does \emph{not} reject $H_0$; this section presents the
forecaster diagnostics that drive that outcome and the design
ablation that confirms the allocator itself is calibrated rather than
mis-specified.

\subsection{Data, Splits and Window Selection}
\label{sec:defi-data-narrative}

The training panel is the $13{,}096$-hour Aave V3 / $12{,}904$-hour
Compound V3 grid of \cref{sec:window}, joined to the
$12{,}895$-hour $98.5\%$-coverage overlap of \cref{sec:coverage}.
The chronological 70/15/15 split of \cref{sec:splits} places the
binding test fold in the final four months of the panel
(2026-01 -- 2026-04). This window was not chosen retrospectively: it
straddles the 2025-Q3-to-Q4 cross-protocol regime inversion documented
in \cref{sec:empirical-spread,sec:regime-empirical} and continues
into the partial 2026-Q2 Aave-dominant phase. It is the most
adversarial single test fold the panel admits, and was selected
deliberately to test transferability under distribution shift rather
than under in-sample stationarity. A more lenient validation-split
choice (e.g.\ a 2025-Q1 test fold inside the Compound-calm regime)
would inflate the headline numbers without testing the architectural
claim of \cref{sec:lob-bridge}.

\subsection{Forecaster Training and Diagnostics}
\label{sec:defi-forecaster-narrative}

The DA-BiGRU-CNN of \cref{sec:dual-branch} is trained with the
composite loss of \cref{eq:loss} under the
$(\alpha,\beta,\gamma)=(0.4,0.5,0.1)$ schedule transferred verbatim
from the LOB specialisation of \cite{solovev2026lob}. The selected
configuration (\cref{sec:fcst-grid}) has hidden dimension $64$, two
BiGRU layers per branch, multi-scale 1-D CNN kernels
$\{3,5,7\}$, dropout $0.1$, learning rate $10^{-3}$, batch size $32$,
sequence length $T=168$ hours and forecast horizon $H=12$ hours; the
exported ONNX module is $1.26$\,MB. Training was conducted on a
Kaggle P100 instance with the v9 hyperparameter pin; the recorded
best-validation loss is $0.427$ and elapsed wall-clock is
$21.5$ minutes for the converged run.

On the validation fold the network attains weighted-Pearson $0.747$
on Aave and $0.289$ on Compound, with directional accuracy $0.739$
on the binary ``Aave $>$ Compound at $t+12\mathrm{h}$'' label and a
positive in-sample $R^{2}=0.138$ on the Compound branch. The
Compound $R^{2}$ exceeds the constant-mean predictor, confirming that
the dual-branch design extracts non-trivial signal on a protocol whose
base rate must be reconstructed via RPC view-call replays
(\cref{sec:data-sources}). The Aave-side $R^{2}$ is negative on the
raw rate level ($-2.12$) but this is the expected signature of the
kink-subtracted residual design: the dominant determinist component
sits in $f_{\text{kink}}(u)$ and the architecture learns the residual
$\varepsilon = r-f_{\text{kink}}(u)$, so absolute rate-level $R^{2}$
under-counts the architecturally-extracted signal that the
weighted-Pearson score correctly weighs.

On the held-out test fold the picture inverts. Weighted-Pearson is
$-0.07$ on Aave and $+0.60$ on Compound; directional accuracy on the
cross-protocol label collapses to $0.35$, below the $0.5$ random-coin
baseline. The pooled out-of-sample diagnostics reported in
\cref{tab:forecast-quality} are dominated by this collapse: the
$\ForecastDirAcc{}$ directional accuracy is below the random baseline
and the pooled rate-level $R^{2}=\ForecastRsq$ is deeply negative
once the kink-residual is reconstituted. The 2025-Q3-to-Q4 regime
crossover documented in \cref{sec:regime-empirical} --- the
Aave-pays-more share rising from $39.9\%$ to $56.3\%$, with median
spread inverting from $-0.218$\,pp to $+0.183$\,pp --- is the most
plausible attribution. The validation fold sits inside the
pre-inversion regime; the test fold straddles the inversion and a
return to Compound dominance in 2026-Q1. The sub-chance directional
accuracy is the empirical signature of a forecaster that has
high-confidence in well-calibrated weighted-Pearson terms locked
onto a correlation pattern from the wrong regime: the model is not
indecisive on the test fold, it commits to the wrong side of the
cross-protocol comparison.

\subsection{Headline H1 Test}
\label{sec:defi-h1-narrative}

The forecast vector feeds the 4-factor MCDM allocator of
\cref{sec:mcdm,eq:mcdm} (APY, Risk, Cost, Stability with default
weights $(0.40,0.25,0.20,0.15)$), with rebalance gated by the
three-gate protocol of \cref{alg:rebalance}: a hysteresis threshold
$\theta=0.05$ on the criterion-score delta, a cooldown
$\tau=1$\,hour, and a gas-cost gate requiring expected uplift over
the forecast horizon to exceed realized gas spend. The reactive
comparator is the EMA-smoothed MCDM-EMA of \cite{solovev2026vault}
re-instantiated as a sibling subclass of
\texttt{BaseLendingAllocationStrategy} on the same backtest harness;
both strategies execute on identical observation streams to ensure
apples-to-apples comparison.

The binding numbers, reproduced from \cref{tab:baselines} for
self-containment, are: predictive net APY \PredictiveAPY{} against
reactive \EMAAPY{}, an $8$\,bp annualized gap in favor of the
reactive baseline; predictive Sharpe \PredictiveSharpe{} against
reactive \EMASharpe{}; both strategies rebalance twice on the
window, with gas spend \PredictiveGas{} versus \EMAGas{} dollars
respectively. The cumulative equity curves (figure logged under
\texttt{results/figures/main\_vs\_ema\_equity.png}) track each
other to within ten basis points over the four-month window; the
predictive variant pays its overhead but does not earn it back.

The $1000$-resample paired monthly-Sharpe bootstrap (predictive
minus reactive on per-month Sharpe pairs) returns
$\Delta\textsc{Sharpe}=\SharpeDelta{}$ with $95\%$ confidence
interval $[\BootstrapCILow{},\,\BootstrapCIHigh{}]$ and bootstrap
$p$-value $\BootstrapPvalue{}$. Both endpoints of the CI lie below
the $H_1$ threshold of $+0.2$; the verdict is therefore
\emph{fail to reject $H_0$}. The wide CI magnitude is a
four-month-window artifact (only four paired monthly observations
enter the resample) and is the same caveat documented in the
pre-registration; with a longer test fold the same point estimate
would carry a tighter interval, but the sign of the contrast would
not change. The bootstrap distribution is concentrated below zero
on every resample, which is why $p=\BootstrapPvalue{}$ for the
one-sided $H_1$ alternative: the resampling procedure cannot find a
sub-window in which the predictive variant outperforms by the
$+0.2$ Sharpe pre-registration threshold.

\subsection{Per-Quarter Regime Decomposition}
\label{sec:defi-regime}

The test window splits cleanly along the two empirically distinct
regimes documented in \cref{sec:regime-empirical}: 2026-Q1
(Compound-calm, $27.1\%$ Aave-pays-more share,
$\sigma=0.57$\,pp) and the partial 2026-Q2 (Aave-dominant,
$61.3\%$ Aave-pays-more share, $\sigma=3.61$\,pp). The
regime-conditional breakdown of \cref{tab:regime-results} reports
the two strategies separately on each sub-window. In Q1 the
predictive variant Sharpe is \QOneSharpePred{} versus EMA
\QOneSharpeEMA{}, with predictive net APY \QOneAPYPred{} versus
EMA \QOneAPYEMA{}; the gap is $1.1$\,pp annualized in favor of the
reactive comparator on this quarter alone, attributable to
unnecessary rebalances triggered by a forecaster that has not yet
recognized the regime as Compound-dominant. In Q2 the predictive
and reactive strategies converge to identical numbers
(Sharpe \QTwoSharpePred{}, APY \QTwoAPYPred{}) because both
implementations elect to remain in the dominant venue throughout the
sub-window, and the forecaster's directional error does not cost
gas. The Q1-vs-Q2 contrast confirms that the H1 negative result is
not driven by gas overhead in steady state: it is driven by
forecaster mis-direction at the regime boundary, where the cost of
the wrong call is highest.

\subsection{Hysteresis-Sweep Ablation}
\label{sec:defi-ablation-narrative}

A natural rebuttal to the H1 negative result is that the default
hysteresis threshold $\theta=0.05$ is mis-calibrated and that the
predictive variant could outperform under a different rebalance
cadence. Ablation~\#16 sweeps $\theta\in\{0.01, 0.02, 0.05, 0.10,
0.20\}$ on the predictive MCDM with all other parameters fixed
(\cref{tab:ablations} row 16 and the row C/E entries in the
ablation macro block). Increasing rebalance frequency from $0$
($\theta\in\{0.10, 0.20\}$) to $4$ ($\theta\in\{0.01, 0.02\}$)
\emph{hurts} risk-adjusted return: net APY moves from
$\AblationSixteenDAPY{}$ at zero rebalances (no action triggered) to
$\AblationSixteenAAPY{}$ at four rebalances, even though Sharpe sits
in the $\AblationSixteenASharpe{}$--$\AblationSixteenCSharpe{}$
range. The default $\theta=0.05$ (row C) sits at the knee of the
curve, with two rebalances, $\AblationSixteenCSharpe{}$ Sharpe, and
$\AblationSixteenCAPY{}$ APY --- gas-light and risk-adjusted, exactly
the calibration the pre-registration of \cref{sec:strat-grid}
prescribed. The ablation does not change the H1 verdict, but it
does establish that the design is correctly tuned: the gap to the
EMA baseline is not a hysteresis-knob artifact, and no setting of
$\theta$ within the swept range recovers a positive
$\Delta\textsc{Sharpe}$ against the reactive comparator.

\subsection{Summary}
\label{sec:defi-summary}

In summary: the architecture extracts in-sample structure on both
protocols (validation weighted-Pearson $0.747$ on Aave with
directional accuracy $0.739$, a positive Compound branch
$R^{2}=0.138$); the allocator is correctly calibrated (default
$\theta=0.05$ sits at the Sharpe knee of the hysteresis sweep, with
two rebalances and gas-light execution); and yet the headline $H_1$
fails on the binding test fold, with
$\Delta\textsc{Sharpe}=\SharpeDelta{}$, $95\%$ CI
$[\BootstrapCILow{},\,\BootstrapCIHigh{}]$ and
$p=\BootstrapPvalue{}$. The proximate cause is the
quarterly-time-scale regime structure documented in
\cref{sec:regime-empirical} that the four-month test window slices
adversarially. The discussion of \cref{sec:discussion} unpacks what
this means for the LOB-to-DeFi cross-domain transfer hypothesis at
each of the three layers --- architectural, predictive, and
strategic.
